% Primary and supplementary source basis: :contentReference[oaicite:3]{index=3} :contentReference[oaicite:4]{index=4}

\section{Polar Form}

\begin{conceptbox}
    If
    \[
        z=x+iy,
    \]
    then we may also write it in polar form as
    \[
        z=r(\cos\theta+i\sin\theta)=re^{i\theta},
    \]
    where
    \[
        r=|z|=\sqrt{x^2+y^2},
        \qquad
        x=r\cos\theta,
        \qquad
        y=r\sin\theta.
    \]
    Here $r$ is the \textbf{modulus} and $\theta$ is an \textbf{argument} of $z$.
\end{conceptbox}

Polar form is especially useful because multiplication, division, powers, and roots become much simpler than in rectangular form.

\subsection*{Motivation / Intuition}

In rectangular form, multiplication mixes real and imaginary parts. In polar form, multiplication becomes:

\[
    r_1e^{i\theta_1}\cdot r_2e^{i\theta_2}
    =
    (r_1r_2)e^{i(\theta_1+\theta_2)}.
\]

So multiplication means:

\begin{itemize}
    \item multiply the magnitudes,
    \item add the angles.
\end{itemize}

That is exactly why polar form is powerful in engineering applications involving gain and phase.

\subsection*{Geometric Interpretation}

The number $z=re^{i\theta}$ is the point whose distance from the origin is $r$ and whose angle from the positive real axis is $\theta$.

Because angles differing by multiples of $2\pi$ represent the same direction,
\[
    \theta,\quad \theta+2\pi,\quad \theta+4\pi,\dots
\]
all correspond to the same complex number.

\begin{conceptbox}
    The argument is \textbf{not unique}. If $\theta$ is one argument of $z\neq 0$, then every argument is
    \[
        \theta+2k\pi, \qquad k\in\mathbb{Z}.
    \]
    A commonly chosen value is the \textbf{principal argument}, denoted $\mathrm{Arg}\,z$.
\end{conceptbox}

\subsection*{Key Properties / Theorems}

\begin{conceptbox}
    Euler's formula:
    \[
        e^{i\theta}=\cos\theta+i\sin\theta.
    \]

    De Moivre's theorem:
    \[
        (re^{i\theta})^n = r^n e^{in\theta}
        = r^n\big(\cos(n\theta)+i\sin(n\theta)\big).
    \]
\end{conceptbox}

\begin{conceptbox}
    For nonzero complex numbers,
    \[
        z_1=r_1e^{i\theta_1}, \qquad z_2=r_2e^{i\theta_2},
    \]
    we have
    \[
        z_1z_2=r_1r_2e^{i(\theta_1+\theta_2)},
    \]
    \[
        \frac{z_1}{z_2}=\frac{r_1}{r_2}e^{i(\theta_1-\theta_2)}.
    \]
\end{conceptbox}

\begin{examplebox}
    \textbf{Example 1: Convert to polar form}

    Express
    \[
        z=-\sqrt{6}-i\sqrt{2}
    \]
    in polar form.

    First compute the modulus:
    \[
        r=\sqrt{(-\sqrt{6})^2+(-\sqrt{2})^2}
        =
        \sqrt{6+2}
        =
        2\sqrt{2}.
    \]

    Now determine the argument. Since
    \[
        \tan\theta=\frac{-\sqrt{2}}{-\sqrt{6}}=\frac{1}{\sqrt{3}},
    \]
    the reference angle is
    \[
        \frac{\pi}{6}.
    \]

    Because $z$ lies in the third quadrant, the correct argument is
    \[
        \theta=\pi+\frac{\pi}{6}=\frac{7\pi}{6}.
    \]

    Hence,
    \[
        z=2\sqrt{2}\,e^{i7\pi/6}.
    \]
\end{examplebox}

\begin{examplebox}
    \textbf{Example 2: Compute a power using De Moivre's theorem}

    Find
    \[
        (1+i)^4.
    \]

    First convert $1+i$ to polar form.

    Its modulus is
    \[
        r=\sqrt{1^2+1^2}=\sqrt{2},
    \]
    and its argument is
    \[
        \theta=\frac{\pi}{4}.
    \]

    Thus,
    \[
        1+i=\sqrt{2}\,e^{i\pi/4}.
    \]

    Raise to the fourth power:
    \[
        (1+i)^4 = \left(\sqrt{2}\,e^{i\pi/4}\right)^4
        = (\sqrt{2})^4 e^{i\pi}
        =4(-1)=-4.
    \]
\end{examplebox}

\begin{examplebox}
    \textbf{Example 3: Find cube roots}

    Find all cube roots of
    \[
        -1+i.
    \]

    Write the number in polar form. Its modulus is
    \[
        r=\sqrt{(-1)^2+1^2}=\sqrt{2},
    \]
    and one argument is
    \[
        \phi=\frac{3\pi}{4}.
    \]

    So
    \[
        -1+i=\sqrt{2}\,e^{i3\pi/4}.
    \]

    The cube roots are
    \[
        z_k = r^{1/3} e^{i(\phi+2k\pi)/3},
        \qquad k=0,1,2.
    \]

    Since
    \[
        r^{1/3}=(\sqrt{2})^{1/3}=2^{1/6},
    \]
    we obtain
    \[
        z_k = 2^{1/6} e^{i\left(\frac{3\pi}{4}+2k\pi\right)/3},
        \qquad k=0,1,2.
    \]

    Thus:
    \[
        z_0=2^{1/6}e^{i\pi/4},
        \qquad
        z_1=2^{1/6}e^{i11\pi/12},
        \qquad
        z_2=2^{1/6}e^{i19\pi/12}.
    \]
\end{examplebox}

\begin{noteBox}
    In frequency-response analysis, a transfer function value is often written as
    \[
        H(i\omega)=|H(i\omega)|e^{i\phi(\omega)}.
    \]
    The modulus gives gain and the argument gives phase shift. This is one reason polar form is central in electrical and control engineering.
\end{noteBox}

\subsection*{Common Mistakes / Notes}

\begin{itemize}
    \item Do not determine $\theta$ using only $\tan^{-1}(y/x)$ without checking the quadrant.
    \item Remember that the argument is multivalued:
          \[
              \arg z = \theta + 2k\pi.
          \]
    \item For roots, do not stop at one answer. An $n$th root problem has $n$ distinct roots.
\end{itemize}